% !PRINTED 508
\weberpage{508}

To this is joined the \gesperrt{theorem of Lagrange}\footnote{\gesperrt{Lagrange}, Réflexions sur la résolution algébrique des équations. Mémoires de l'Académie de Berlin, années 1770, 1771. Oeuvres de Lagrange. Tome III. The theorem was, to be sure, given by \gesperrt{Lagrange} only in a more special form. The general formulation is due to \gesperrt{Galois}.}:

2. \gesperrt{Every quantity of the field $N$ which admits the permutations
 of the group $Q$ is contained in the
field $\Omega(\psi)$, if $\psi$ is a function belonging
 to $Q$.}

A function $\omega$ which admits the permutations of $Q$ passes,
under the permutations of a subsidiary group $Q\pi_1$, into one and
the same function. Hence to the conjugate
values

\weq{3}{\psi, \psi_1, \psi_2 \ldots \psi_{j-1}}

there correspond the values

\weq{4}{\omega, \omega_1, \omega_2 \ldots \omega_{j-1},}

which, however, are not necessarily all different from one another.

If one applies to the sequences of quantities (3), (4) one of the
permutations of $P$, a certain permutation takes place, and
indeed the same one in both sequences, since, if for example $\psi_1$ passes into $\psi_2$,
then $\omega_1$ must also pass into $\omega_2$ (§. 154, 5.).

Let us now consider the sum

\[
\varphi(t) \left( \frac{\omega}{t - \psi} + \frac{\omega_1}{t - \psi_1} + \cdots + \frac{\omega_{j-1}}{t - \psi_{j-1}} \right) = \chi(t),
\]

which is an integral function of $t$ of degree $(j-1)$; we then find
that it remains unchanged under all the permutations $P$ and
is consequently contained in $\Omega$. If one then sets $t = \psi$ and observes
that $\varphi(t)$ has no equal roots, it follows that

\weq{5}{\omega = \frac{\chi(\psi)}{\varphi'(\psi)}.}

Thus $\omega$ is expressed rationally in terms of $\psi$, and this is the
content of theorem 2.

3. \gesperrt{If we adjoin to the field $\Omega$ a function $\psi$ belonging
to $Q$, the group of the
field $N$ reduces to $Q$.}

% [page alignment mirrored 2026-08-07] Printed page 508 contains the complete
% line through "because $\psi$ remains".  Page 509 begins "unchanged by $Q$";
% the English boundary now follows the source-confirmed German recut.
For if we denote the field $\Omega(\psi)$ by $\Omega'$, then \gesperrt{first},
every equation in $\Omega'$ between $\alpha, \alpha_1 \ldots \alpha_{m-1}$ admits the permutations
of $Q$, because $Q$ is contained in $P$, and because $\psi$ remains
%CONT>>
